\chapter{Notatki terenowe}\label{dod:notatki}

\begin{blok}
Zapisuj tu obserwacje z terenu, poprawki do procedur i pomysły. Każdą notatkę opatrz datą, nazwą misji i podpisem. Po misji przenieś notatki do repozytorium dokumentacji (formularz F14).
\end{blok}

% Strony w linie
\newcommand{\stronanotatek}{%
  \clearpage
  \noindent{\sffamily\small\color{sdgray}Misja: \rule{45mm}{0.3pt}\hfill Data: \rule{30mm}{0.3pt}\hfill Autor: \rule{40mm}{0.3pt}}\par
  \vspace{2mm}
  \noindent\begin{tikzpicture}
    \foreach \y in {0,8,...,224}{\draw[sdgray!45,line width=0.3pt] (0,-\y mm) -- (\linewidth,-\y mm);}
  \end{tikzpicture}}
% Strony w kratkę (szkice)
\newcommand{\stronakratka}{%
  \clearpage
  \noindent{\sffamily\small\color{sdgray}Szkic -- misja: \rule{40mm}{0.3pt}\hfill Data: \rule{30mm}{0.3pt}\hfill Autor: \rule{40mm}{0.3pt}}\par
  \vspace{2mm}
  \noindent\begin{tikzpicture}
    \draw[sdgray!30,line width=0.2pt,step=5mm] (0,0) grid (\linewidth,225mm);
  \end{tikzpicture}}

\stronanotatek
\stronanotatek
\stronanotatek
\stronanotatek
\stronakratka
\stronakratka
